\documentclass[a4j]{jsarticle}

\begin{document}

\title{はじめての\LaTeX}
\author{2022311042 神野友哉}
\date{\today}%2022年6月1日}
\maketitle

\section{はじめに}
…
\subsection{研究目的}
…
\section{準備}
…
\subsection{数学的定義}
…

\subsection{問題の記述}
…
\begin{itemize}
  \item ファイルシステム
  \item Webページ
  \item 文書整形
\end{itemize}
\begin{enumerate}
  \item C言語
  \item Javascript
\end{enumerate}
こんにちは, \verb5\LaTeX5 .
\\今回が人生初の \LaTeX によるレポート作成です.


私は情報科学部の神野友哉です.
\\今後ともよろしくお願いします.


Hello. I`m Tomoya Shinno.
\\Nice meeting you.
\section{表の例題}
  \begin{tabular}{l|c|r}
\hline
\hline
品名　&　個数 & 価格 \\
\hline
燕の子安貝　& 2 & 11000 \\
龍の首の珠　& 1 & 5000 \\
蓬莱の玉の枝 & 3 & 8500 \\
\hline
\end{tabular}

\section{数式の例}
\begin{equation} 
 \left|
  \begin{array}{ccc}
      a & b & c \\
      d & e & f \\
      g & h & j \\
    \end{array}
  \right| = aei + bfg + cdh - ( afh + bdi + ceg )
\end{equation}
\section{導関数}
 関数 $f(x)$　の導関数 $f'(x)$は、
  \begin{equation}
    f'(x)=\lim_{h\rightarrow 0}\frac{f(x+h)-f(x)}{h}
  \end{equation}
 で定義される。例えば、関数 $f(x)=3x^2+2x$ の導関数は、
 \[ f'(x)=6x+2 \]
である。
\end{document}
